% 数理物理2 数理物理2-1 問題3（003-demoivre） 解いた日 2026-11-04
%% kotae
\toi{1}{◎}{}
（先にオイラーの公式で $\cos3\theta$・$\sin3\theta$ を求めている。見出しは「(2)」と書いてある）
\begin{align*}
\cos3\theta&=\frac{1}{2}(\ee^{3\ii\theta}+\ee^{-3\ii\theta})\\
&=\frac{1}{2}\bigl((\ee^{\ii\theta}+\ee^{-\ii\theta})(\ee^{2\ii\theta}-1+\ee^{-2\ii\theta})\bigr)\\
&=\cos\theta\bigl((\ee^{\ii\theta}+\ee^{-\ii\theta})^{2}-3\bigr)\\
&=\kakimachigai{\cos}{\cos\theta}(4\cos^{2}\theta-3)\\
&=4\cos^{3}\theta-3\cos\theta
\end{align*}
\begin{align*}
\sin3\theta&=\frac{1}{2\ii}(\ee^{3\ii\theta}-\ee^{-3\ii\theta})\\
&=\frac{1}{2\ii}\bigl((\ee^{\ii\theta}-\ee^{-\ii\theta})(\ee^{2\ii\theta}+1+\ee^{-2\ii\theta})\bigr)\\
&=\sin\theta\bigl((\ee^{\ii\theta}+\ee^{-\ii\theta})^{2}-1\bigr)\\
&=\sin\theta\bigl(4(1-\sin^{2}\theta)-1\bigr)\\
&=3\sin\theta-4\sin^{3}\theta
\end{align*}
（続けて見出し「(1)」でド・モアブルの定理を使う）
\begin{align*}
(\cos\theta+\ii\sin\theta)^{3}&=\cos3\theta+\ii\sin3\theta\\
&=\cos^{3}\theta+3\ii\cos^{2}\theta\sin\theta\\
&\qquad+3\cos\theta(\ii\sin\theta)^{2}+(\ii\sin\theta)^{3}\\
&=\cos^{3}\theta-3\cos\theta(1-\cos^{2}\theta)\\
&\qquad+\ii\bigl(3\sin\theta(1-\sin^{2}\theta)-\sin^{3}\theta\bigr)\\
&=4\cos^{3}\theta-3\cos\theta+\ii(3\sin\theta-4\sin^{3}\theta)
\end{align*}
\begin{align*}
\therefore\ \cos3\theta&=4\cos^{3}\theta-3\cos\theta\\
\sin3\theta&=3\sin\theta-4\sin^{3}\theta
\end{align*}
\toi{2}{◎}{}
\begin{align*}
\cos^{3}\theta&=\Bigl(\frac{\ee^{\ii\theta}+\ee^{-\ii\theta}}{2}\Bigr)^{3}\\
&=\frac{1}{8}(\ee^{3\ii\theta}+3\ee^{\ii\theta}+3\ee^{-\ii\theta}+\ee^{-3\ii\theta})\\
&=\frac{1}{4}\cos3\theta+\frac{3}{4}\cos\theta
\end{align*}
\begin{align*}
\cos3\theta&=4\cos^{3}\theta-3\cos\theta\\
&=(\cos3\theta+3\cos\theta)-3\cos\theta\\
&=\cos3\theta
\end{align*}
より一致。
\toi{3}{◎}{}
\begin{align*}
\sin^{4}\theta&=\Bigl(\frac{\ee^{\ii\theta}-\ee^{-\ii\theta}}{\kakimachigai{2}{2\ii}}\Bigr)^{4}\\
&=\frac{1}{16}(\ee^{4\ii\theta}-4\ee^{2\ii\theta}+6-4\ee^{-2\ii\theta}+\ee^{-4\ii\theta})\\
&=\frac{1}{8}\cos4\theta-\frac{1}{2}\cos2\theta+\frac{3}{8}
\end{align*}
\begin{align*}
I&=\frac{1}{8}\int_{0}^{2\pi}\kakimachigai{\cos4\theta-4\cos2\theta+3}{(\cdots)}\,\dd\theta\\
&=\frac{1}{8}\Bigl[\frac{1}{4}\sin4\theta-2\sin2\theta+3\theta\Bigr]_{0}^{2\pi}\\
&=\frac{1}{8}(6\pi-0)=\frac{3}{4}\pi
\end{align*}
%% kizuki
\item (3) で $\sin\theta=(\ee^{\ii\theta}-\ee^{-\ii\theta})/(2\ii)$ の分母を $2$ と書いた。$(2\ii)^{4}=16=2^{4}$ なので4乗では結果が変わらないが，奇数乗では符号と $\ii$ が残る。$\sin$ の分母は $2\ii$ と覚え，$\theta=\pi/2$ を入れて $\sin\theta=1$ になるかで確かめる。
\item 被積分関数が2項以上のときは括弧でくくってから $\dd\theta$ を書く（$\int(\cos4\theta-4\cos2\theta+3)\,\dd\theta$）。
